\documentclass[11pt, a4paper]{article}

% ── 패키지 ────────────────────────────────────────────────────
\usepackage{fontspec}
\usepackage{geometry}
\usepackage{xcolor}
\usepackage{titlesec}
\usepackage{titletoc}
\usepackage{fancyhdr}
\usepackage{hyperref}
\usepackage{listings}
\usepackage{booktabs}
\usepackage{tabularx}
\usepackage{array}
\usepackage{longtable}
\usepackage{enumitem}
\usepackage{parskip}
\usepackage{microtype}
\usepackage{mdframed}
\usepackage{graphicx}
\usepackage{amsmath}
\usepackage{setspace}
\usepackage{pifont}
\usepackage{tcolorbox}
\tcbuselibrary{listings, skins, breakable}

% ── 페이지 레이아웃 ───────────────────────────────────────────
\geometry{
  a4paper,
  top=25mm, bottom=28mm,
  left=25mm, right=25mm,
  headheight=14pt,
  headsep=10mm,
  footskip=12mm
}

% ── 폰트 ──────────────────────────────────────────────────────
\setmainfont{TeX Gyre Termes}
\setmonofont{DejaVu Sans Mono}[Scale=0.85]

% ── 색상 ──────────────────────────────────────────────────────
\definecolor{accent}{RGB}{37, 99, 235}       % 파란 액센트
\definecolor{accentlight}{RGB}{219,234,254}  % 연한 파랑 배경
\definecolor{codebg}{RGB}{248, 249, 250}     % 코드 배경
\definecolor{codefg}{RGB}{36, 41, 47}        % 코드 텍스트
\definecolor{codekw}{RGB}{215, 58, 73}       % 키워드
\definecolor{codestr}{RGB}{3, 136, 47}       % 문자열
\definecolor{codecmt}{RGB}{106, 115, 125}    % 주석
\definecolor{borderclr}{RGB}{208, 215, 222}  % 테두리
\definecolor{headerbg}{RGB}{246, 248, 250}   % 헤더 배경
\definecolor{ruleblue}{RGB}{37, 99, 235}
\definecolor{notebg}{RGB}{254, 252, 232}
\definecolor{noteborder}{RGB}{202, 138, 4}
\definecolor{warnbg}{RGB}{254, 242, 242}
\definecolor{warnborder}{RGB}{220, 38, 38}

% ── 하이퍼링크 ────────────────────────────────────────────────
\hypersetup{
  colorlinks=true,
  linkcolor=accent,
  urlcolor=accent,
  citecolor=accent,
  pdftitle={PCS Analyzer Plugin Distribution Guide},
  pdfauthor={PCS Analyzer},
  pdfsubject={Plugin Distribution Guide v0.1},
}

% ── 섹션 스타일 ───────────────────────────────────────────────
\titleformat{\section}
  {\Large\bfseries\color{black}}
  {\color{accent}\thesection}{1em}{}
  [\vspace{1pt}\color{accent}\rule{\linewidth}{1.5pt}\vspace{-4pt}]

\titleformat{\subsection}
  {\normalsize\bfseries\color{black}}
  {\color{accent}\thesubsection}{1em}{}

\titleformat{\subsubsection}
  {\normalsize\bfseries\itshape\color{black}}
  {\thesubsubsection}{1em}{}

\titlespacing{\section}{0pt}{18pt}{8pt}
\titlespacing{\subsection}{0pt}{12pt}{4pt}
\titlespacing{\subsubsection}{0pt}{8pt}{3pt}

% ── 목차 스타일 ───────────────────────────────────────────────
\titlecontents{section}[0em]
  {\smallskip\bfseries}
  {\contentslabel{2.2em}}{}
  {\titlerule*[0.5pc]{.}\contentspage}

\titlecontents{subsection}[2.2em]
  {\small}
  {\contentslabel{2.8em}}{}
  {\titlerule*[0.5pc]{.}\contentspage}

% ── 헤더/푸터 ─────────────────────────────────────────────────
\pagestyle{fancy}
\fancyhf{}
\fancyhead[L]{\small\color{gray}\textit{PCS Analyzer Plugin Distribution Guide}}
\fancyhead[R]{\small\color{gray}v0.1 — Draft}
\fancyfoot[C]{\small\color{gray}\thepage}
\renewcommand{\headrulewidth}{0.4pt}
\renewcommand{\footrulewidth}{0pt}
\renewcommand{\headrule}{\color{borderclr}\hrule width\headwidth height\headrulewidth}

% ── 코드 블록 ─────────────────────────────────────────────────
\lstdefinestyle{python}{
  language=Python,
  backgroundcolor=\color{codebg},
  basicstyle=\ttfamily\small\color{codefg},
  keywordstyle=\bfseries\color{codekw},
  stringstyle=\color{codestr},
  commentstyle=\itshape\color{codecmt},
  numberstyle=\tiny\color{gray},
  numbers=none,
  breaklines=true,
  breakatwhitespace=false,
  showspaces=false,
  showstringspaces=false,
  tabsize=4,
  frame=single,
  framesep=6pt,
  rulecolor=\color{borderclr},
  xleftmargin=8pt,
  xrightmargin=8pt,
  aboveskip=6pt,
  belowskip=6pt,
  literate={~}{{\textasciitilde}}{1},
  morekeywords={self, True, False, None},
}

\lstdefinestyle{json}{
  language=,
  backgroundcolor=\color{codebg},
  basicstyle=\ttfamily\small\color{codefg},
  stringstyle=\color{codestr},
  breaklines=true,
  frame=single,
  framesep=6pt,
  rulecolor=\color{borderclr},
  xleftmargin=8pt,
  xrightmargin=8pt,
  aboveskip=6pt,
  belowskip=6pt,
  morestring=[b]",
}

\lstdefinestyle{text}{
  language=,
  backgroundcolor=\color{codebg},
  basicstyle=\ttfamily\small\color{codefg},
  breaklines=true,
  frame=single,
  framesep=6pt,
  rulecolor=\color{borderclr},
  xleftmargin=8pt,
  xrightmargin=8pt,
  aboveskip=6pt,
  belowskip=6pt,
}

% ── 노트 / 경고 박스 ──────────────────────────────────────────
\newmdenv[
  backgroundcolor=notebg,
  linecolor=noteborder,
  linewidth=1pt,
  leftline=true, rightline=false, topline=false, bottomline=false,
  leftmargin=0pt, rightmargin=0pt,
  innerleftmargin=10pt, innerrightmargin=8pt,
  innertopmargin=6pt, innerbottommargin=6pt,
  skipabove=8pt, skipbelow=8pt,
]{notebox}

\newmdenv[
  backgroundcolor=warnbg,
  linecolor=warnborder,
  linewidth=1pt,
  leftline=true, rightline=false, topline=false, bottomline=false,
  leftmargin=0pt, rightmargin=0pt,
  innerleftmargin=10pt, innerrightmargin=8pt,
  innertopmargin=6pt, innerbottommargin=6pt,
  skipabove=8pt, skipbelow=8pt,
]{warnbox}

% ── 테이블 스타일 ─────────────────────────────────────────────
\newcolumntype{L}{>{\raggedright\arraybackslash}X}
\newcolumntype{M}[1]{>{\raggedright\arraybackslash}p{#1}}

\setlength{\arrayrulewidth}{0.4pt}
\arrayrulecolor{borderclr}

% ── 리스트 ────────────────────────────────────────────────────
\setlist[itemize]{
  leftmargin=1.5em,
  itemsep=2pt,
  parsep=0pt,
  topsep=4pt,
}
\setlist[enumerate]{
  leftmargin=1.8em,
  itemsep=2pt,
  parsep=0pt,
  topsep=4pt,
}

% ── 체크리스트 ────────────────────────────────────────────────
\newlist{checklist}{itemize}{1}
\setlist[checklist]{
  label={\color{borderclr}\ding{111}},
  leftmargin=1.6em,
  itemsep=1pt,
  parsep=0pt,
  topsep=4pt,
}

% ══════════════════════════════════════════════════════════════
\begin{document}

% ── 커버 페이지 ───────────────────────────────────────────────
\begin{titlepage}
  \newgeometry{left=25mm, right=25mm, top=40mm, bottom=30mm}

  % 상단 액센트 바
  {\color{accent}\rule{\linewidth}{4pt}}
  \vspace{10mm}

  {\fontsize{9}{11}\selectfont\color{gray}\textls[80]{PCS ANALYZER}}\\[4mm]
  {\fontsize{28}{34}\bfseries\selectfont Plugin Distribution Guide}\\[6mm]
  {\color{accent}\rule{60mm}{1.5pt}}\\[6mm]
  {\large\color{gray} Version 0.1 \quad|\quad Draft Specification}

  \vspace{10mm}
  {\normalsize
    \begin{tabular}{@{}ll}
      \textbf{Document status:} & Draft specification \\[2pt]
      \textbf{Target application:} & PCS Analyzer \\[2pt]
      \textbf{Scope:} & External plugin authors, internal module developers, \\
                      & and future maintainers \\
    \end{tabular}
  }

  \vfill

  {\normalsize\color{gray}
    This document defines the recommended packaging, installation, loading,\\
    and distribution rules for PCS Analyzer plugins.
  }

  \vspace{16mm}
  {\color{accent}\rule{\linewidth}{1pt}}
  \restoregeometry
\end{titlepage}

% ── 목차 ──────────────────────────────────────────────────────
\clearpage
\tableofcontents
\clearpage

% ══════════════════════════════════════════════════════════════
\section{Core Design Principle}

PCS Analyzer plugins should be as independent as possible.

The main application provides only:

\begin{itemize}
  \item the \texttt{Modules} menu
  \item the plugin loading system
  \item the Module Manager
  \item a minimal \texttt{PluginApp} API
  \item optional read-only access to selected PCS Analyzer state
\end{itemize}

A plugin should not require modification of any core PCS Analyzer source file, including \texttt{ui/components.py}, \texttt{logic/}, \texttt{ui/}, or \texttt{main.py}.

A plugin registers itself through a \texttt{register(app)} function. The main application loads the plugin and calls:

\begin{lstlisting}[style=python]
plugin_module.register(plugin_app)
\end{lstlisting}

The plugin is responsible for adding its own menu entry through the provided \texttt{PluginApp} object.

PCS Analyzer supports two user-facing plugin distribution styles:
\begin{enumerate}
  \item \textbf{Single-file Python plugins} (\texttt{.py})
  \item \textbf{Folder-based plugins}, optionally distributed as \texttt{.zip} packages
\end{enumerate}

Internally, all plugins are installed as folder-based packages under:
\begin{lstlisting}[style=text]
plugins/installed/<plugin_id>/
\end{lstlisting}

Even when a user installs a single \texttt{.py} plugin, PCS Analyzer converts it into the standard installed folder layout during installation.

% ──────────────────────────────────────────────────────────────
\section{Terminology}

\begin{tabularx}{\linewidth}{@{}M{3.6cm}L@{}}
  \toprule
  \textbf{Term} & \textbf{Meaning} \\
  \midrule
  Plugin & An external module installed under \texttt{plugins/installed/<plugin\_id>/}. \\
  Built-in module & A tool shipped as part of the PCS Analyzer core distribution, such as PCS Workbench. \\
  Module Manager & GUI used to install, enable, disable, and remove external plugins. \\
  PluginApp & Minimal API wrapper passed to plugins by PCS Analyzer. \\
  Registry & The \texttt{plugins/plugins.json} file that stores installed plugin metadata and enabled state. \\
  Entry file & The Python file loaded by the plugin loader, usually \texttt{plugin.py}. \\
  \bottomrule
\end{tabularx}

\vspace{4pt}
Both built-in modules and external plugins may be listed under the \texttt{Modules} menu. The Module Manager controls only external plugins.

% ──────────────────────────────────────────────────────────────
\section{Standard Installed Plugin Layout}

Every installed plugin must have the following internal structure:

\begin{lstlisting}[style=text]
plugins/installed/<plugin_id>/
├─ manifest.json
└─ plugin.py
\end{lstlisting}

For complex plugins, additional folders may be included:

\begin{lstlisting}[style=text]
plugins/installed/<plugin_id>/
├─ manifest.json
├─ plugin.py
├─ plugin_ui/
│  ├─ __init__.py
│  └─ main_window.py
├─ plugin_logic/
│  ├─ __init__.py
│  ├─ models.py
│  └─ fitting.py
├─ resources/
│  └─ icon.png
└─ examples/
   └─ sample_data.csv
\end{lstlisting}

The plugin loader reads \texttt{manifest.json}, finds the entry file, imports it, and calls \texttt{register(app)}.

% ──────────────────────────────────────────────────────────────
\section{Single-File Plugin Distribution}

Single-file plugins are recommended for small utilities, such as simple calculators, data converters, quick plotting tools, small export helpers, and prototype analysis tools.

A single-file plugin may be distributed as a plain \texttt{.py} file, for example \texttt{evans\_calculator.py}. It must contain at minimum a \texttt{PLUGIN\_INFO} dictionary and a \texttt{register(app)} function:

\begin{lstlisting}[style=python]
PLUGIN_INFO = {
    "id": "evans_calculator",
    "name": "Evans Method Calculator",
    "version": "0.1.0",
    "author": "Author Name",
    "description": "Calculate magnetic susceptibility from Evans method data.",
    "type": "window",
    "standalone": True,
    "dependencies": [],
}


def register(app):
    app.add_menu_item(
        label="Evans Method Calculator...",
        command=lambda: open_plugin(app),
    )


def open_plugin(app):
    ...
\end{lstlisting}

When installed through the Module Manager, PCS Analyzer converts it to a standard folder layout. The generated \texttt{manifest.json} is derived from \texttt{PLUGIN\_INFO}. Users should not need to manually create \texttt{manifest.json} for single-file plugins.

% ──────────────────────────────────────────────────────────────
\section{Folder-Based Plugin Distribution}

Folder-based plugins are recommended for larger analysis modules, such as Bleaney/VT shift fitting, VT-NMR line-shape analysis, CShM/coordination geometry analysis, crystal-field fitting utilities, and PCS ensemble or motion analysis tools.

A folder-based plugin should be distributed as a folder containing:

\begin{enumerate}
  \item a valid \texttt{manifest.json}
  \item the entry file specified by \texttt{manifest.json}, usually \texttt{plugin.py}
  \item a \texttt{register(app)} function in the entry file
\end{enumerate}

% ──────────────────────────────────────────────────────────────
\section{ZIP Plugin Distribution}

For public distribution, complex plugins may be packaged as \texttt{.zip} files.

Recommended structure:

\begin{lstlisting}[style=text]
bleaney_vt.zip
└─ bleaney_vt/
   ├─ manifest.json
   ├─ plugin.py
   ├─ plugin_ui/
   ├─ plugin_logic/
   ├─ resources/
   └─ examples/
\end{lstlisting}

During installation, PCS Analyzer should:

\begin{enumerate}
  \item Extract the zip to a temporary folder
  \item Locate \texttt{manifest.json}
  \item Validate the plugin ID and entry file
  \item Copy the plugin folder to \texttt{plugins/installed/<plugin\_id>/}
  \item Register the plugin in \texttt{plugins/plugins.json}
  \item Refresh the \texttt{Modules} menu if supported
\end{enumerate}

% ──────────────────────────────────────────────────────────────
\section{\texttt{manifest.json} Format}

Every installed folder-based plugin must contain a \texttt{manifest.json}. Below is a minimum example followed by a recommended full example.

\subsection{Minimum Example}

\begin{lstlisting}[style=json]
{
  "id": "test_plugin",
  "name": "Test Plugin",
  "version": "0.1.0",
  "author": "Author Name",
  "description": "Simple plugin loading test.",
  "entry": "plugin.py",
  "type": "window",
  "standalone": true
}
\end{lstlisting}

\subsection{Recommended Full Example}

\begin{lstlisting}[style=json]
{
  "id": "bleaney_vt",
  "name": "Bleaney / VT Shift Fitting",
  "version": "0.1.0",
  "author": "Author Name",
  "description": "Temperature-dependent paramagnetic shift fitting module.",
  "entry": "plugin.py",
  "type": "window",
  "standalone": true,
  "min_app_version": "1.4.0",
  "dependencies": ["numpy", "scipy", "pandas", "matplotlib"],
  "optional_dependencies": ["pyvista", "pyvistaqt", "imageio"],
  "state_access": "read_only",
  "category": "Paramagnetic NMR"
}
\end{lstlisting}

\subsection{Required Fields}

\begin{tabularx}{\linewidth}{@{}M{3.2cm}L@{}}
  \toprule
  \textbf{Field} & \textbf{Description} \\
  \midrule
  \texttt{id} & Stable and unique plugin identifier. Use lowercase letters, numbers, and underscores. \\
  \texttt{name} & Display name shown in the Module Manager. \\
  \texttt{version} & Plugin version. \\
  \texttt{entry} & Entry Python file, usually \texttt{plugin.py}. \\
  \bottomrule
\end{tabularx}

\subsection{Recommended Optional Fields}

\begin{tabularx}{\linewidth}{@{}M{3.8cm}L@{}}
  \toprule
  \textbf{Field} & \textbf{Description} \\
  \midrule
  \texttt{author} & Plugin author. \\
  \texttt{description} & Short plugin description. \\
  \texttt{type} & Plugin type: \texttt{window}, \texttt{analysis}, \texttt{viewer}, \texttt{export}, or \texttt{tool}. \\
  \texttt{standalone} & Whether the plugin can run independently. \\
  \texttt{min\_app\_version} & Minimum PCS Analyzer version required. \\
  \texttt{dependencies} & Required Python packages. \\
  \texttt{optional\_dependencies} & Optional packages used for enhanced functionality. \\
  \texttt{state\_access} & Intended access level: \texttt{none}, \texttt{read\_only}, or \texttt{read\_write}. \\
  \texttt{category} & Display or grouping category. \\
  \bottomrule
\end{tabularx}

\subsection{Deprecated Aliases}

New plugins should use \texttt{"dependencies"} and \texttt{"optional\_dependencies"}. The following older aliases may be accepted by the installer for compatibility but should not be used in new plugins: \texttt{"requires"} and \texttt{"optional\_requires"}.

\subsection{Reserved Fields}

The field \texttt{"entry\_function"} is reserved for possible future loader support. Current PCS Analyzer plugin loading is based on \texttt{register(app)}. Plugins must not rely on \texttt{entry\_function} unless the loader is explicitly updated to support it.

% ──────────────────────────────────────────────────────────────
\section{Plugin ID Rules}

The plugin ID must be stable and unique. Use lowercase letters, numbers, and underscores.

\textbf{Good examples:}
\begin{lstlisting}[style=text]
bleaney_vt    evans_calculator    condon_helper
vt_lineshape  coordination_geometry_analyzer
\end{lstlisting}

\textbf{Avoid:}
\begin{lstlisting}[style=text]
Bleaney Module    my plugin!!!    test    plugin1
\end{lstlisting}

The plugin ID is used as the installation folder name, registry key, import namespace, enable/disable state identifier, and update/reinstall matching key. Changing the plugin ID should be treated as creating a new plugin.

% ──────────────────────────────────────────────────────────────
\section{Entry File Rules}

The entry file is usually \texttt{plugin.py}. It must define:

\begin{lstlisting}[style=python]
def register(app):
    ...
\end{lstlisting}

For a single-file plugin, relative imports should be replaced by local functions or absolute imports. For a folder-based plugin, see Section~\ref{sec:imports}.

% ──────────────────────────────────────────────────────────────
\section{Import-Time Safety Rules}

\begin{warnbox}
\textbf{Important:} This is one of the most critical rules. A plugin must not perform any side effects at import time.
\end{warnbox}

A plugin must not perform any of the following at import time:

\begin{itemize}
  \item open a GUI window
  \item call \texttt{main()} or \texttt{mainloop()}
  \item start calculations or launch subprocesses
  \item load large files unnecessarily
  \item modify PCS Analyzer state
  \item show message boxes (unless dependency validation explicitly requires it)
\end{itemize}

Only definitions, constants, lightweight imports, and small helper functions should run at import time.

\subsection*{Incorrect}
\begin{lstlisting}[style=python]
from my_tool.standalone import main
main()  # launches immediately on import
\end{lstlisting}

\subsection*{Correct}
\begin{lstlisting}[style=python]
def open_plugin(app):
    from my_tool.standalone import main
    return main()

def register(app):
    app.add_menu_item(
        label="My Tool...",
        command=lambda: open_plugin(app),
    )
\end{lstlisting}

% ──────────────────────────────────────────────────────────────
\section{Menu Callback Rules}

Plugin menu callbacks must be deferred using \texttt{lambda}.

\subsection*{Correct}
\begin{lstlisting}[style=python]
app.add_menu_item(
    label="My Plugin...",
    command=lambda: open_plugin(app),
)
\end{lstlisting}

\subsection*{Incorrect}
\begin{lstlisting}[style=python]
app.add_menu_item(
    label="My Plugin...",
    command=open_plugin(app),  # called immediately during registration
)
\end{lstlisting}

% ──────────────────────────────────────────────────────────────
\section{Folder Plugin Import Rules}
\label{sec:imports}

Folder-based plugins should use package-style relative imports:

\begin{lstlisting}[style=python]
from .plugin_ui.main_window import MyWindow
from .plugin_logic.fitting import run_fit
\end{lstlisting}

This requires the plugin loader to use \texttt{submodule\_search\_locations} and register the module in \texttt{sys.modules} before execution:

\begin{lstlisting}[style=python]
spec = importlib.util.spec_from_file_location(
    module_name, entry_path,
    submodule_search_locations=[str(plugin_dir)],
)
module = importlib.util.module_from_spec(spec)
sys.modules[module_name] = module
spec.loader.exec_module(module)
\end{lstlisting}

If package-style loading is unavailable, local absolute imports may be used, but this is less robust and may conflict with other plugins.

% ──────────────────────────────────────────────────────────────
\section{Standalone Compatibility}

A plugin may also run as a standalone tool. The standalone entry point must be protected:

\begin{lstlisting}[style=python]
def run_standalone():
    import tkinter as tk
    root = tk.Tk()
    root.title("My Plugin")
    MyWindow(root)
    root.mainloop()

if __name__ == "__main__":
    run_standalone()
\end{lstlisting}

Never call standalone code directly at import time.

% ──────────────────────────────────────────────────────────────
\section{\texttt{PluginApp} API Usage}

Plugins should interact with PCS Analyzer through the \texttt{PluginApp} wrapper. Common API methods include:

\begin{lstlisting}[style=python]
app.root                          # parent Tk window
app.state                         # application state dict (read-only)
app.add_menu_item(label, command) # register menu entry
app.add_separator()               # add menu separator
app.get_current_structure()
app.get_delta_exp_values()
app.get_pcs_values_by_id()
app.get_metal_position()
app.refresh_views()
\end{lstlisting}

Direct modification of \texttt{app.state} should be avoided unless the plugin is explicitly designed as an integrated plugin.

% ──────────────────────────────────────────────────────────────
\section{State Access Levels}

\begin{tabularx}{\linewidth}{@{}M{2.8cm}L@{}}
  \toprule
  \textbf{Value} & \textbf{Meaning} \\
  \midrule
  \texttt{none} & Plugin does not use PCS Analyzer state. \\
  \texttt{read\_only} & Plugin reads structures, shifts, PCS values, or settings. \\
  \texttt{read\_write} & Plugin may modify state. Use only for tightly integrated plugins. \\
  \bottomrule
\end{tabularx}

\vspace{4pt}
Most plugins should declare \texttt{"state\_access": "none"} or \texttt{"state\_access": "read\_only"}.

% ──────────────────────────────────────────────────────────────
\section{Dependency Handling}

List required dependencies in \texttt{manifest.json}:

\begin{lstlisting}[style=json]
"dependencies": ["numpy", "scipy", "matplotlib"],
"optional_dependencies": ["pyvista", "pyvistaqt", "imageio"]
\end{lstlisting}

Plugins should handle missing optional dependencies gracefully and check for required dependencies at runtime without launching the GUI at import time.

% ──────────────────────────────────────────────────────────────
\section{GUI Integration Rules}

\begin{itemize}
  \item Use \texttt{tk.Toplevel(app.root)} for child windows when running inside PCS Analyzer.
  \item Do not call \texttt{tk.Tk()} or \texttt{mainloop()} when running inside PCS Analyzer.
  \item \texttt{tk.Tk()} and \texttt{mainloop()} are allowed only in standalone mode.
\end{itemize}

% ──────────────────────────────────────────────────────────────
\section{Security Rules}

Python plugins can execute arbitrary code on the user's computer. The Module Manager should display an appropriate warning and require confirmation before installing any plugin.

Plugin authors should avoid: unexpected network access, modifying user files without confirmation, hidden subprocess execution, modifying PCS Analyzer core files, and storing credentials or private data.

% ──────────────────────────────────────────────────────────────
\section{Update and Reinstallation Policy}

Updating a plugin may require restarting PCS Analyzer. On Windows, an imported plugin may prevent its folder from being deleted.

Recommended update strategy:

\begin{enumerate}
  \item Rename the existing folder to a backup (e.g., \texttt{bleaney\_vt.\_\_old\_\_20260516/}).
  \item Copy the new plugin folder into \texttt{plugins/installed/<plugin\_id>/}.
  \item Update \texttt{plugins/plugins.json}.
  \item Refresh the \texttt{Modules} menu.
  \item Attempt to delete the backup folder.
  \item If deletion fails, leave a \texttt{.delete\_on\_restart} marker and notify the user.
\end{enumerate}

% ──────────────────────────────────────────────────────────────
\section{Module Manager Behavior}

The Module Manager should support: listing, enabling, disabling, removing, and installing plugins (\texttt{.py}, folder, or \texttt{.zip}).

\begin{notebox}
\textbf{Note:} Visible in Module Manager $\neq$ successfully loaded into the main menu.
\end{notebox}

If a plugin appears in the Module Manager but not in the main menu, likely causes include: \texttt{enabled} is false, \texttt{manifest.json} is missing or invalid, \texttt{register(app)} is missing, an exception during import or registration, unsupported relative imports, or missing required dependencies.

% ──────────────────────────────────────────────────────────────
\section{Registry Format: \texttt{plugins/plugins.json}}

\begin{lstlisting}[style=json]
{
  "vt_nmr_analysis": {
    "id": "vt_nmr_analysis",
    "name": "VT-NMR Analysis",
    "version": "3.2.0",
    "dir": "vt_nmr_analysis",
    "enabled": true
  }
}
\end{lstlisting}

The registry should be treated as application-managed data. Plugin authors should not ask users to edit it manually except for debugging.

% ──────────────────────────────────────────────────────────────
\section{Loader Requirements}

A robust PCS Analyzer plugin loader should:

\begin{enumerate}
  \item Read \texttt{plugins/plugins.json} and skip disabled plugins
  \item Resolve \texttt{plugins/installed/<plugin\_id>/} and read \texttt{manifest.json}
  \item Validate the entry file
  \item Import using package-style loading and register in \texttt{sys.modules}
  \item Verify and call \texttt{register(app)}
  \item Store loading results and errors
  \item Avoid crashing the whole application when one plugin fails
\end{enumerate}

\begin{lstlisting}[style=python]
module_name = f"pcs_plugin_{plugin_id}"
spec = importlib.util.spec_from_file_location(
    module_name, entry_path,
    submodule_search_locations=[str(plugin_dir)],
)
module = importlib.util.module_from_spec(spec)
sys.modules[module_name] = module
spec.loader.exec_module(module)
if not hasattr(module, "register"):
    raise RuntimeError("register(app) function is missing.")
module.register(app)
\end{lstlisting}

% ──────────────────────────────────────────────────────────────
\section{Recommended Plugin Folder Naming}

Use plugin-local folder names to avoid conflicts with PCS Analyzer core modules.

\textbf{Recommended:} \texttt{plugin\_ui/}, \texttt{plugin\_logic/}, \texttt{resources/}, \texttt{examples/}

\textbf{Avoid:} \texttt{ui/}, \texttt{logic/}, \texttt{core/}, \texttt{utils/}

% ──────────────────────────────────────────────────────────────
\section{Resource and Data File Access}

Plugins should resolve resource paths relative to the plugin file:

\begin{lstlisting}[style=python]
from pathlib import Path

PLUGIN_DIR   = Path(__file__).resolve().parent
RESOURCE_DIR = PLUGIN_DIR / "resources"
EXAMPLE_DIR  = PLUGIN_DIR / "examples"

# Correct
open(RESOURCE_DIR / "icon.png", "rb")

# Incorrect - do not rely on working directory
open("resources/icon.png")
\end{lstlisting}

% ──────────────────────────────────────────────────────────────
\section{Versioning Policy}

Plugin versions should follow semantic versioning:

\begin{tabularx}{\linewidth}{@{}M{3cm}L@{}}
  \toprule
  \textbf{Part} & \textbf{Meaning} \\
  \midrule
  \texttt{MAJOR} & Breaking changes. \\
  \texttt{MINOR} & New features without breaking compatibility. \\
  \texttt{PATCH} & Bug fixes. \\
  \bottomrule
\end{tabularx}

% ──────────────────────────────────────────────────────────────
\section{Compatibility Policy}

Plugins may declare the minimum required PCS Analyzer version via \texttt{"min\_app\_version"}. Future plugin managers may refuse to load plugins requiring a newer version. Plugins should avoid relying on undocumented internal state keys unless declared as \texttt{state\_access: "read\_write"}.

% ──────────────────────────────────────────────────────────────
\section{Logging and Error Handling}

Plugins should catch user-facing errors and report them clearly:

\begin{lstlisting}[style=python]
try:
    run_analysis()
except Exception as exc:
    from tkinter import messagebox
    messagebox.showerror("My Plugin", f"Analysis failed:\n\n{exc}")
\end{lstlisting}

Plugin loaders should catch import and registration errors and continue loading other plugins.

% ──────────────────────────────────────────────────────────────
\section{Distribution Checklist}

Before distributing a plugin, verify all of the following:

\begin{checklist}
  \item The plugin has a stable lowercase plugin ID.
  \item The plugin defines \texttt{PLUGIN\_INFO} or \texttt{manifest.json}.
  \item Folder-based plugins include \texttt{manifest.json}.
  \item The manifest has \texttt{id}, \texttt{name}, \texttt{version}, and \texttt{entry}.
  \item The entry file defines \texttt{register(app)}.
  \item The plugin does not open windows at import time.
  \item The plugin does not call \texttt{main()} at import time.
  \item The plugin does not call \texttt{mainloop()} inside PCS Analyzer.
  \item Menu callbacks use \texttt{command=lambda: open\_plugin(app)}.
  \item Required dependencies are listed.
  \item Missing optional dependencies are handled gracefully.
  \item Resource paths are resolved relative to the plugin folder.
  \item The plugin does not modify PCS Analyzer core source files.
  \item The plugin can be removed without breaking the core application.
  \item The plugin was tested after installation through Module Manager.
  \item The plugin was tested after disabling and re-enabling.
  \item The plugin was tested after reinstalling/updating.
\end{checklist}

% ──────────────────────────────────────────────────────────────
\section{Minimal Single-File Plugin Template}

\begin{lstlisting}[style=python]
from __future__ import annotations
import tkinter as tk
from tkinter import ttk

PLUGIN_INFO = {
    "id": "minimal_plugin",
    "name": "Minimal Plugin",
    "version": "0.1.0",
    "author": "Author Name",
    "description": "A minimal PCS Analyzer plugin.",
    "type": "window",
    "standalone": True,
    "dependencies": [],
}

def open_plugin(app):
    win = tk.Toplevel(app.root)
    win.title("Minimal Plugin")
    win.geometry("360x160")
    frame = ttk.Frame(win, padding=16)
    frame.pack(fill="both", expand=True)
    ttk.Label(frame, text="Minimal plugin loaded successfully.").pack(anchor="w")
    return win

def register(app):
    app.add_menu_item(
        label="Minimal Plugin...",
        command=lambda: open_plugin(app),
    )

def run_standalone():
    root = tk.Tk()
    root.title("Minimal Plugin")
    class DummyApp:
        root = root
        state = {}
    open_plugin(DummyApp())
    root.mainloop()

if __name__ == "__main__":
    run_standalone()
\end{lstlisting}

% ──────────────────────────────────────────────────────────────
\section{Minimal Folder-Based Plugin Template}

\subsection{Folder Structure}

\begin{lstlisting}[style=text]
my_analysis_plugin/
├─ manifest.json
├─ plugin.py
├─ plugin_ui/
│  ├─ __init__.py
│  └─ main_window.py
└─ plugin_logic/
   ├─ __init__.py
   └─ analysis.py
\end{lstlisting}

\subsection{\texttt{manifest.json}}

\begin{lstlisting}[style=json]
{
  "id": "my_analysis_plugin",
  "name": "My Analysis Plugin",
  "version": "0.1.0",
  "author": "Author Name",
  "description": "Example folder-based analysis plugin.",
  "entry": "plugin.py",
  "type": "analysis",
  "standalone": true,
  "dependencies": ["numpy", "matplotlib"],
  "state_access": "read_only"
}
\end{lstlisting}

\subsection{\texttt{plugin.py}}

\begin{lstlisting}[style=python]
from __future__ import annotations

PLUGIN_INFO = {
    "id": "my_analysis_plugin",
    "name": "My Analysis Plugin",
    "version": "0.1.0",
}

def open_plugin(app):
    from .plugin_ui.main_window import MyAnalysisWindow
    return MyAnalysisWindow(parent=app.root, app=app)

def register(app):
    app.add_menu_item(
        label="My Analysis Plugin...",
        command=lambda: open_plugin(app),
    )

def run_standalone():
    import tkinter as tk
    from .plugin_ui.main_window import MyAnalysisWindow
    root = tk.Tk()
    root.title("My Analysis Plugin")
    MyAnalysisWindow(parent=root, app=None)
    root.mainloop()

if __name__ == "__main__":
    run_standalone()
\end{lstlisting}

% ──────────────────────────────────────────────────────────────
\section{Common Failure Modes}

\subsection{Plugin appears in Module Manager but not in the main menu}

Likely causes:

\begin{itemize}
  \item \texttt{enabled} is \texttt{false} in the registry
  \item \texttt{register(app)} is missing
  \item Import failed due to missing dependency or relative import error
  \item Plugin opened itself and crashed during import
\end{itemize}

Check \texttt{state["plugin\_load\_results"]} or console diagnostics.

\subsection{Plugin opens automatically on startup}

The menu callback was written as \texttt{command=open\_plugin(app)} instead of \texttt{command=lambda: open\_plugin(app)}, or \texttt{main()} was called at import time.

\subsection{Relative import fails}

\begin{lstlisting}[style=text]
ImportError: attempted relative import with no known parent package
\end{lstlisting}

Fix: update the loader to use \texttt{submodule\_search\_locations}, or avoid relative imports and use local absolute imports. Package-style relative imports are preferred.

\subsection{Windows refuses to update/delete plugin folder}

\begin{lstlisting}[style=text]
PermissionError: [WinError 5] Access is denied
\end{lstlisting}

Likely causes: plugin is currently imported, plugin window is open, or the OS is indexing files. Close plugin windows, restart PCS Analyzer, and try again. The installer should use backup-folder replacement instead of direct deletion.

% ──────────────────────────────────────────────────────────────
\section{Recommended Development Workflow}

\begin{enumerate}
  \item Develop as standalone tool first.
  \item Add \texttt{PLUGIN\_INFO}.
  \item Add \texttt{open\_plugin(app)}.
  \item Add \texttt{register(app)}.
  \item Ensure no code runs at import time.
  \item Install through Module Manager.
  \item Test loading from the \texttt{Modules} menu.
  \item Test disable/enable.
  \item Test reinstall/update.
  \item Package as \texttt{.py}, folder, or \texttt{.zip}.
\end{enumerate}

% ──────────────────────────────────────────────────────────────
\section{Summary of Mandatory Rules}

A valid PCS Analyzer plugin must obey the following rules:

\begin{enumerate}
  \item It must have a stable plugin ID.
  \item It must define \texttt{register(app)}.
  \item It must not execute the tool at import time.
  \item It must use deferred menu callbacks: \texttt{command=lambda: open\_plugin(app)}.
  \item It must not modify PCS Analyzer core files.
  \item It must declare dependencies when distributed publicly.
  \item It must handle missing optional dependencies gracefully.
  \item It must keep standalone execution behind \texttt{if \_\_name\_\_ == "\_\_main\_\_"}.
\end{enumerate}

These rules are intended to keep PCS Analyzer modular, stable, and safe as the plugin ecosystem grows.

\end{document}
